%% chapters/01-pendahuluan.tex
%% Bab I — Pendahuluan
%% Aturan: lihat .cursor/rules/06-bab1-pendahuluan.mdc

\chapter{Pendahuluan}
\label{bab:pendahuluan}

\section{Latar Belakang}
\label{subsec:bab1_latar}

Bantalan gelinding (\emph{rolling element bearings}) adalah komponen
mekanis yang sangat luas digunakan dalam mesin rotari modern, mulai
dari motor listrik berkecepatan tinggi hingga turbin angin
berkapasitas \si{\mega\watt}. Kegagalan bantalan menjadi penyebab
sekitar XX\,\% dari seluruh kerusakan mesin rotari di industri
\citetitb{Reference1}, sehingga kemampuan memprediksi sisa umur pakai
(\emph{remaining useful life}, RUL) bantalan secara akurat menjadi
fondasi strategi pemeliharaan prediktif (\emph{predictive maintenance}).

% [paragraph 2: tren PHM dan deep learning]
% [paragraph 3: keterbatasan SOTA — LSTM saturasi pada sekuens panjang;
%  Transformer mahal komputasi; minim interpretabilitas]
% [paragraph 4: peluang Mamba dan xLSTM, motivasi hibrida]
% [paragraph 5: pentingnya interpretabilitas — sparse autoencoder]
% [paragraph 6: posisi dan ringkasan kontribusi disertasi ini]

\section{Masalah Penelitian}
\label{subsec:bab1_masalah}

Berdasarkan latar belakang di atas, dirumuskan tiga pertanyaan
penelitian (\emph{research questions}, RQ) sebagai berikut:

\begin{description}
  \item[RQ1.] Bagaimana mendesain arsitektur \emph{deep learning} yang
        memodelkan dinamika sinyal getaran jangka panjang sekaligus
        mempertahankan presisi rekuren untuk prediksi RUL bantalan
        yang lebih akurat dibanding pendekatan SOTA berbasis
        Transformer maupun LSTM klasik?
  \item[RQ2.] Bagaimana modul interpretabilitas \emph{post-hoc} dapat
        dirancang sehingga representasi laten model RUL dapat
        dipetakan secara empiris ke konsep fisis yang dapat
        diverifikasi dengan teori klasik diagnosis getaran bantalan?
  \item[RQ3.] Sejauh mana arsitektur dan metodologi yang diusulkan
        dapat digeneralisasi lintas dataset dengan kondisi operasi
        yang berbeda?
\end{description}

\section{Tujuan Penelitian}
\label{subsec:bab1_tujuan}

Penelitian ini bertujuan untuk:

\begin{enumerate}
  \item Mengembangkan arsitektur hibrida \emph{Mamba-xLSTM-Net} yang
        menggabungkan keunggulan \emph{state-space model} selektif
        dan \emph{matrix-memory recurrent} untuk prediksi RUL
        bantalan gelinding.
  \item Merancang dan memvalidasi modul interpretabilitas berbasis
        \emph{Top-k Sparse Autoencoder} yang dapat dipetakan ke
        frekuensi karakteristik bantalan.
  \item Memvalidasi metodologi yang diusulkan secara empiris pada
        dataset publik PHM2012 dan XJTU-SY dengan baseline
        komprehensif.
\end{enumerate}

\section{Lingkup Penelitian}
\label{subsec:bab1_lingkup}

Penelitian ini dibatasi pada:

\begin{itemize}
  \item Bantalan gelinding jenis \emph{deep groove ball bearing} pada
        rezim degradasi alami (\emph{natural run-to-failure}).
  \item Sinyal getaran satu sumbu (vertikal) dengan frekuensi sampling
        \SI{25.6}{\kilo\hertz}.
  \item Dataset publik PHM2012 (FEMTO-PRONOSTIA) dan XJTU-SY.
  \item Prediksi RUL skalar (regresi tunggal); tidak mencakup
        klasifikasi mode kegagalan.
  \item Interpretabilitas pada level fitur laten dan korespondensi
        spektral; tidak mencakup interpretabilitas kausal melalui
        intervensi.
\end{itemize}

\section{Asumsi Penelitian}
\label{subsec:bab1_asumsi}

Asumsi yang digunakan:

\begin{enumerate}
  \item Degradasi bantalan diasumsikan monoton secara umum, dengan
        toleransi terhadap fluktuasi lokal.
  \item Label RUL dihitung secara linier menurun dari titik awal
        prediksi (\emph{first prediction time}, FPT) hingga \emph{end
        of life} (EOL).
  \item Karakteristik bantalan dan kondisi operasi pada setiap
        dataset dianggap homogen sesuai dokumentasi resmi dataset.
\end{enumerate}

\section{Hipotesis Penelitian}
\label{subsec:bab1_hipotesis}

\begin{description}
  \item[H1.] Kombinasi \emph{state-space model} selektif (Mamba) dan
        \emph{matrix-memory recurrent} (mLSTM) memberikan kinerja
        prediksi RUL yang lebih baik dibanding masing-masing komponen
        secara terpisah.
  \item[H2.] Modul \emph{Top-k Sparse Autoencoder} yang diterapkan
        secara \emph{post-hoc} dapat menghasilkan fitur laten yang
        konsisten dengan frekuensi karakteristik bantalan tanpa
        mendegradasi akurasi prediksi RUL.
  \item[H3.] Arsitektur yang diusulkan menunjukkan generalisasi yang
        lebih baik lintas dataset dibanding baseline berbasis
        Transformer.
\end{description}

\section{Pendekatan Penelitian}
\label{subsec:bab1_pendekatan}

Penelitian ini menggunakan pendekatan eksperimental kuantitatif
dengan urutan: studi literatur, perumusan arsitektur, implementasi
prototipe, eksperimen sistematis dengan ablasi, analisis
interpretabilitas, dan validasi lintas dataset. Detail metodologi
dijabarkan pada Bab~\ref{bab:metodologi}.

\section{Kebaruan Penelitian}
\label{subsec:bab1_kebaruan}

Disertasi ini menyumbangkan tiga kebaruan utama (kategorisasi
sesuai Pedoman ITB §V.1):

\begin{enumerate}
  \item \textbf{(Teknologi-Metodologi)} Arsitektur hibrida
        \emph{Mamba-xLSTM-Net} yang mengintegrasikan
        \emph{state-space model} selektif dengan
        \emph{matrix-memory xLSTM} untuk prediksi RUL bantalan,
        dengan mekanisme fusi yang dirancang berdasarkan sifat
        non-stasioner sinyal getaran.
  \item \textbf{(Teknologi-Metodologi)} Skema interpretabilitas
        \emph{Top-k Sparse Autoencoder} yang bekerja secara
        \emph{post-hoc} terhadap representasi laten model RUL,
        dilengkapi prosedur kuantitatif untuk memetakan fitur SAE
        ke frekuensi karakteristik fisik bantalan.
  \item \textbf{(Output)} Bukti empiris bahwa model \emph{deep
        learning} yang dilatih untuk RUL mempelajari representasi
        laten yang berkorespondensi dengan teori klasik diagnosis
        getaran bantalan, disertai peningkatan kinerja SOTA pada
        dua dataset publik.
\end{enumerate}

\section{Sistematika Penulisan}
\label{subsec:bab1_sistematika}

Disertasi ini disusun ke dalam enam bab. \textbf{Bab~I} memuat latar
belakang, masalah, tujuan, lingkup, asumsi, hipotesis, pendekatan,
kebaruan, dan sistematika penulisan. \textbf{Bab~II} menyajikan
tinjauan pustaka secara kronologis terkait metode prediksi RUL
bantalan dan perkembangan arsitektur \emph{deep learning}.
\textbf{Bab~III} memuat dasar teori meliputi formulasi RUL, LSTM,
xLSTM (sLSTM dan mLSTM), \emph{state-space model} dan Mamba,
\emph{Sparse Autoencoder}, serta metode atribusi.
\textbf{Bab~IV} menjabarkan metodologi penelitian, termasuk
arsitektur yang diusulkan, dataset, protokol pelatihan, dan
rancangan eksperimen. \textbf{Bab~V} menyajikan hasil eksperimen,
studi ablasi, dan analisis interpretabilitas. \textbf{Bab~VI}
menutup disertasi dengan kesimpulan dan saran pekerjaan lanjut.
